\documentclass[11pt,a4paper]{article}
\usepackage[T1]{fontenc}
\usepackage[utf8]{inputenc}
\usepackage[main=italian,provide=*]{babel}
\usepackage{amsmath,amssymb}
\usepackage{geometry}
\geometry{margin=2.3cm,top=2.5cm,bottom=2.7cm}
\usepackage{booktabs}
\usepackage[table]{xcolor}
\usepackage{tcolorbox}
\tcbuselibrary{skins,breakable}
\usepackage{titlesec}
\usepackage{enumitem}
\usepackage{needspace}
\usepackage{fancyhdr}
\usepackage{hyperref}
\hypersetup{colorlinks=true,linkcolor=Sepia,citecolor=Sepia,urlcolor=Sepia}

% --- palette ---
\definecolor{inkblue}{HTML}{1B3A5C}
\definecolor{lightblue}{HTML}{EAF1F8}
\definecolor{accent}{HTML}{B0562A}
\definecolor{softgray}{HTML}{6B6B6B}
\definecolor{okgreen}{HTML}{2E6B3E}
\definecolor{okgreenbg}{HTML}{EAF5EC}
\definecolor{advisorpurple}{HTML}{5B3A7A}
\definecolor{advisorbg}{HTML}{F1ECF6}

% --- titoli di sezione ---
\titleformat{\section}
  {\normalfont\Large\bfseries\color{inkblue}}
  {\thesection}{0.6em}{}[{\color{inkblue!40}\titlerule}]
\titlespacing{\section}{0pt}{0.7em}{0.4em}
\titleformat{\subsection}
  {\normalfont\large\bfseries\color{accent}}
  {}{0em}{}
\titlespacing{\subsection}{0pt}{0.5em}{0.2em}

\pagestyle{fancy}
\fancyhf{}
\renewcommand{\headrulewidth}{0.4pt}
\fancyhead[L]{\small\color{softgray}Risultati VQE --- punto di lavoro test 2}
\fancyhead[R]{\small\color{softgray}\thepage}
\fancyfoot[C]{\small\color{softgray}Tesi triennale, Samuele Schiavi --- Relatore: Prof.\ Alessandro Chiesa}

% --- box riutilizzabili ---
\newtcolorbox{keyeq}[1][]{
  colback=lightblue, colframe=inkblue!70, boxrule=0.6pt,
  arc=2pt, left=8pt, right=8pt, top=3pt, bottom=3pt, #1
}
\newtcolorbox{panelbox}[1]{
  colback=white, colframe=softgray!50, boxrule=0.5pt,
  arc=2pt, left=7pt, right=7pt, top=2pt, bottom=2pt,
  fonttitle=\bfseries\color{inkblue}, title={#1}
}
\newtcolorbox{okbox}[1]{
  colback=okgreenbg, colframe=okgreen!60, boxrule=0.6pt,
  arc=2pt, left=7pt, right=7pt, top=2pt, bottom=2pt,
  fonttitle=\bfseries\color{okgreen}, title={#1}
}
\newtcolorbox{advisorbox}[1]{
  colback=advisorbg, colframe=advisorpurple!60, boxrule=0.6pt,
  arc=2pt, left=7pt, right=7pt, top=3pt, bottom=3pt,
  fonttitle=\bfseries\color{advisorpurple}, title={#1}
}

\title{\vspace{-1.5em}\color{inkblue}Ground state VQE --- punto di lavoro test 2\\[0.2em]
\Large\normalfont\color{softgray}$b/J=0.35$,\ $D/J=0.80$,\ $J=1$}
\author{Note di lavoro --- Tesi triennale, Samuele Schiavi\\ Relatore: Prof.\ Alessandro Chiesa}
\date{}

\begin{document}
\sloppy
\maketitle
\thispagestyle{fancy}
\vspace{-2.2em}

File: \texttt{vqe\_test2.py}. Punto di lavoro: $b/J=0.35$, $D/J=0.80$,
$J=1$ (parametri del test 2 in \texttt{relazione\_test\_parametri.md}).

\section*{Metodologia (nessuna scelta nuova)}
\addcontentsline{toc}{section}{Metodologia}

\begin{itemize}[leftmargin=1.3em,itemsep=0.25em]
\item \textbf{Hamiltoniana:} stessa convenzione di \texttt{dimer\_exact.py}
/ \texttt{confronto\_ansatz\_entangler.ipynb},
$$H=J(X_1X_2+Y_1Y_2+Z_1Z_2)+b(Z_1+Z_2)+D(X_1Z_2-Z_1X_2).$$
\item \textbf{Ansatz:} \textbf{PMA-2q$\cdot$3} (blocco RBS + $R_y$
indipendenti sui due qubit, 3 parametri) --- il riferimento canonico
deciso in \texttt{confronto\_ansatz\_entangler.ipynb}.
\item \textbf{Ottimizzazione:} multistart $R=6$ (COBYLA, reinizializzazione
nel ciclo esterno) seguito da un \textbf{polish L-BFGS-B} dal miglior
punto trovato. Il polish è un'aggiunta rispetto al notebook di confronto:
lì l'obiettivo era il confronto qualitativo fra ansatz su tutto un range
di $b/J$, qui serve uno stato preciso a precisione macchina perché
diventerà l'input del circuito con l'ancilla nello stadio successivo ---
un residuo $\mathcal O(10^{-5})$ in energia (quanto lasciava COBYLA da
solo) propagherebbe un errore sistematico non controllato nelle
correlazioni.
\end{itemize}

\section*{Risultato}
\addcontentsline{toc}{section}{Risultato}

\begin{okbox}{Convergenza a precisione macchina}
\begin{center}
\renewcommand{\arraystretch}{1.25}
\begin{tabular}{@{}lc@{}}
\toprule
\textbf{Quantità} & \textbf{Valore}\\
\midrule
$E_0$ (esatto) & $-3.57321145$\\
$E$ (VQE) & $-3.57321145$\\
$|E_\text{VQE}-E_0|$ & $7.75\times10^{-11}$\\
Fidelity & $1.00000000$ (a precisione macchina)\\
$\langle M_z\rangle$ & $-0.034730$\\
\bottomrule
\end{tabular}
\end{center}
\end{okbox}

\medskip
\noindent Stato (base $|00\rangle,|01\rangle,|10\rangle,|11\rangle$,
reale come atteso --- vedi nota di rigore in
\texttt{confronto\_ansatz\_entangler.ipynb} sul perché il DM in forma
$XZ-ZX$ mantiene $H$ reale):

\begin{keyeq}
\centering
$|\psi_0\rangle \simeq 0.2017\,|00\rangle + 0.6648\,|01\rangle -
0.6648\,|10\rangle + 0.2746\,|11\rangle$
\end{keyeq}

\noindent Parametri ottimali: $(p_0,p_1,p_2) = (3.978569,\ 1.226719,\
1.914748)$ rad.

\section*{Lettura fisica}
\addcontentsline{toc}{section}{Lettura fisica}

\begin{itemize}[leftmargin=1.3em,itemsep=0.3em]
\item $\langle M_z\rangle\simeq-0.035$: il ground state resta vicino al
settore $M=0$ (il campo $b/J=0.35$ è modesto rispetto a $J$), ma non è
puro $M=0$ --- la piccola componente $|11\rangle$ ($0.2746$) e
l'asimmetria fra $|00\rangle$ e $|11\rangle$ vengono dal termine DM
($D/J=0.8$, non trascurabile) che mescola i settori.
\item Il peso dominante è sulla combinazione antisimmetrica
$|01\rangle-|10\rangle$ (verso il singoletto
$|S\rangle=(|01\rangle-|10\rangle)/\sqrt2$): coerente con l'aspettativa
che a $D/J\lesssim1$ e campo piccolo il ground state resti "vicino" al
singoletto del caso $D=0,b=0$, deformato dal DM.
\item Il fatto che $\langle\psi_\text{exact}|\psi_\text{VQE}\rangle$ sia
reale e positivo (non solo $|\cdot|=1$) conferma che non c'è una fase
relativa residua da fissare a mano prima di usare questo stato come
input di un circuito.
\end{itemize}

\section*{Dispersione del multistart}
\addcontentsline{toc}{section}{Dispersione del multistart}

Le 6 energie finali di COBYLA (prima del polish) hanno dispersione
$6.73\times10^{-3}$ --- indica che il paesaggio a 3 parametri ha
comunque minimi locali non banali a questo punto di lavoro (diversamente
da $D=0$), motivo in più per non fidarsi del solo COBYLA e applicare il
polish.

\needspace{10\baselineskip}
\section*{Prossimo passo}
\addcontentsline{toc}{section}{Prossimo passo}

Lo stato $|\psi_0^\text{VQE}\rangle$ (equivalentemente i 3 parametri
ottimali dell'ansatz PMA-2q$\cdot$3) è salvato in
\texttt{ground\_state\_test2.npz} ed è pronto per essere preparato come
stato iniziale nel circuito con l'ancilla per le correlazioni
dinamiche.

\begin{advisorbox}{Risposta del relatore ricevuta}
Nessuna restrizione --- sia auto-correlazione (stesso sito) sia
correlazione fra spin diversi sono di interesse, e vale la pena provare
più componenti ($S_z,S_x,S_y$); alcune combinazioni si annulleranno per
simmetria dell'Hamiltoniana (atteso, non un errore). Prima di
implementare il circuito: controllo classico di quali combinazioni sono
strutturalmente nulle a questo punto di lavoro. Vedi
\texttt{log\_decisioni.md} e \texttt{domande\_relatore.md} per il
dettaglio.
\end{advisorbox}

\end{document}
